\section{Why A* expands fewer vertices}\label{sec:efficiency}

Although the worst-case bounds coincide, A* expands far fewer vertices than
Dijkstra's algorithm in practice. This section characterizes the vertices it
expands and measures the effect on a grid with an obstacle.

\subsection{The vertices that A* expands}\label{sec:expanded}

\begin{proposition}\label{prop:expanded}
Let $G$ be a finite graph with nonnegative edge weights, let $t$ be reachable
from $s$, and let $h$ be a consistent heuristic. Then Algorithm~\ref{alg:astar}
has the following two properties.
\begin{enumerate}[label=(\alph*)]
  \item Every vertex $u$ that it extracts satisfies
    $\delta(s,u) + h(u) \le C^{*}$.
  \item Every vertex $u$ with $\delta(s,u) + h(u) < C^{*}$ is extracted before
    $t$.
\end{enumerate}
\end{proposition}

\begin{proof}
(a) By Lemma~\ref{lem:exact}, a vertex $u$ is extracted with key
$g(u) + h(u) = \delta(s,u) + h(u)$. The destination is extracted last, with key
$\delta(s,t) + h(t) = C^{*}$, and by Lemma~\ref{lem:monotone} no earlier key is
larger.

(b) Suppose that such a vertex $u$ is not in $S$ when $t$ is extracted. Let $P$
be a shortest path from $s$ to $u$, and let $x$ be the first vertex of $P$ that
is not in $S$. The argument of the proof of Lemma~\ref{lem:exact} shows that $x$
is in $Q$ with $g(x) = \delta(s,x)$, and Lemma~\ref{lem:pathmono} gives
$f(x) \le \delta(s,u) + h(u) < C^{*}$. The key of $t$ is $C^{*}$, so $x \neq t$,
and \textsc{Extract-Min} would have returned a vertex with a key below $C^{*}$
instead of $t$. This is a contradiction.
\end{proof}

In words, A* expands only vertices whose best route is estimated to cost at
most $C^{*}$. For the zero heuristic, part~(b) says that Dijkstra's algorithm
expands every $v$ with $\delta(s,v) < C^{*}$ (Remark~\ref{rem:dijkstra}).

\begin{corollary}\label{cor:subset}
Let $h$ be a consistent heuristic. Every vertex $u$ that A* expands but
Dijkstra's algorithm does not expand satisfies $\delta(s,u) = C^{*}$ and
$h(u) = 0$.
\end{corollary}

\begin{proof}
Let $u$ be such a vertex. By part~(a) of Proposition~\ref{prop:expanded},
$\delta(s,u) + h(u) \le C^{*}$, and since $h(u) \ge 0$, we have
$\delta(s,u) \le C^{*}$. By part~(b) applied to the zero heuristic, Dijkstra's
algorithm expands every vertex $v$ with $\delta(s,v) < C^{*}$, so
$\delta(s,u) = C^{*}$. Then part~(a) gives $h(u) \le 0$, that is, $h(u) = 0$.
\end{proof}

Apart from such borderline vertices, A* expands a subset of what Dijkstra's
algorithm expands. The set $\{u : \delta(s,u) + h(u) \le C^{*}\}$ shrinks as $h$
grows, which justifies requirement~(iii) of Section~\ref{sec:h-requirements};
for $h = h^{*}$ it contains only the vertices on shortest paths.

On the running example, $\delta(s,v) + h(v)$ is 7 for \cty{D}, \cty{N},
\cty{B}, \cty{H}, 8 for \cty{S}, and at least 9 for the others, so A* expands
exactly the five cities of Table~\ref{tab:trace-h}.

Geometrically, with the zero heuristic this set is a disk around $s$; with the
Euclidean heuristic it is a thin ellipse around the segment from $s$ to $t$
(Figure~\ref{fig:funnel}).

\begin{figure}[htbp]
  \centering
  \begin{tabular}{@{}c@{\hspace{1.5em}}c@{}}
  {\footnotesize $h \equiv 0$ (Dijkstra): a flat floor} &
  {\footnotesize Euclidean $h$ (A*): a funnel into $t$}\\[2pt]
  \begin{tikzpicture}
    \begin{axis}[land3d, land view={-28}{40}{0.30},
                 colormap name=htiles, point meta min=0, point meta max=3]
      \addplot3[surf, shader=flat corner, draw=white, line width=0.45pt,
                samples=19, samples y=13, domain=0:18, y domain=0:12,
                point meta={
                  and(abs(x+0.5-\cSx)<0.1, abs(y+0.5-\cSy)<0.1) ? 2 :
                  (and(abs(x+0.5-\cGx)<0.1, abs(y+0.5-\cGy)<0.1) ? 3 :
                  (veclen(x+0.5-\cSx, y+0.5-\cSy) <= \cCost ? 1 : 0))}]
        {0};
      \draw[edgpath, line width=1.6pt, line cap=round] (\cSx,\cSy,0) -- (\cGx,\cGy,0);
    \end{axis}
  \end{tikzpicture}
  &
  \begin{tikzpicture}
    \begin{axis}[land3d, land view={-28}{40}{0.30},
                 colormap name=htiles, point meta min=0, point meta max=3]
      \addplot3[surf, shader=flat corner, draw=white, line width=0.45pt,
                z buffer=sort, samples=19, samples y=13, domain=0:18, y domain=0:12,
                point meta={
                  and(abs(x+0.5-\cSx)<0.1, abs(y+0.5-\cSy)<0.1) ? 2 :
                  (and(abs(x+0.5-\cGx)<0.1, abs(y+0.5-\cGy)<0.1) ? 3 :
                  (veclen(x+0.5-\cSx, y+0.5-\cSy)
                   + veclen(x+0.5-\cGx, y+0.5-\cGy) <= \cCost+\cSlack ? 1 : 0))}]
        {\kC*veclen(x-\cGx, y-\cGy)};
      \addplot3[edgpath, line width=1.6pt, line cap=round,
                domain=\cSx:\cGx, samples=21, samples y=1]
        ({x}, {\cSy}, {\kC*sqrt((x-\cGx)^2+0.25)});
    \end{axis}
  \end{tikzpicture}\\
  {\scriptsize the search grows as a disk} & {\scriptsize the search is a thin ellipse from $s$ to $t$}\\[4pt]
  \multicolumn{2}{c}{\legendC}
\end{tabular}

  \caption{A schematic comparison on one grid. With the zero heuristic the
    expanded cells form a disk around $s$; with the Euclidean heuristic they
    form a thin ellipse from $s$ to $t$.}\label{fig:funnel}
\end{figure}

\subsection{Experiment: a grid with a U-shaped wall}\label{sec:grid}

An obstacle tests how A* copes with information that its heuristic does not
contain. The instance is as follows (Figure~\ref{fig:grid3d}).

\begin{description}[font=\normalfont\itshape, leftmargin=1.5em, labelindent=0pt]
    \item[Graph.] The 139 free cells of a $16 \times 10$ grid; moves go up, down,
    left or right and cost 1.
  \item[Wall.] A U of 21 cells opening toward the source: column 12, rows 3--7,
    and rows 3 and 7, columns 4--12.
  \item[Query.] From $s = (1,4)$ to $t = (14,6)$, behind the wall, with the
    Manhattan distance as heuristic.
\end{description}

Ties are broken first in, first out. Table~\ref{tab:grid-compare},
Figure~\ref{fig:grid-bars} and Figure~\ref{fig:grid-panels} give the results.

\begin{figure}[htbp]
  \centering
  \begin{tikzpicture}[cam={-24}{54}{0.52}{0.52}, line cap=round, line join=round]
  \gfloorthree
  \gendsthree
  \foreach \c in {0,...,15}{\node[tag, text=neutralg, font=\tiny] at ({\c+0.5},-0.6,0) {\c};}
  \foreach \r in {0,...,9}{\node[tag, text=neutralg, font=\tiny] at (-0.7,{\r+0.5},0) {\r};}
  \gwallsthree
\end{tikzpicture}

  \caption{The $16 \times 10$ test grid. The raised blocks form the U-shaped
    wall; the source $s$ is green and the destination $t$ is red.}\label{fig:grid3d}
\end{figure}

\begin{table}[htbp]
  \centering
  \caption{The three algorithms on the grid with the U-shaped
    wall.}\label{tab:grid-compare}
  \begin{tabular}{@{}lccc@{}}
    \toprule
     & \textcolor{algdijkstra}{\textbf{Dijkstra's algorithm}}
     & \textcolor{alggreedy}{\textbf{Greedy best-first}}
     & \textcolor{algastar}{\textbf{A* search}}\\
    \midrule
    Cells expanded (of 139 free) & 137 & 48 & 81\\
    Path cost (moves)            & 19 & 35 & 19\\
    Extra moves over the optimum & 0 & 16 & 0\\
    Shortest path                & \textcolor{algastar}{yes}
                                 & \textcolor{cbRed}{no}
                                 & \textcolor{algastar}{yes}\\
    \bottomrule
  \end{tabular}
\end{table}

\begin{figure}[htbp]
  \centering
  \resultbars{Cells expanded}{165}{137}{48}{81}{0}\hfill
  \resultbars{Path cost (moves)}{42}{19}{35}{19}{19}
  \caption{The results of Table~\ref{tab:grid-compare} as bar charts. The dashed
    line marks the optimal cost.}\label{fig:grid-bars}
\end{figure}

\begin{figure}[htbp]
  \centering
  \renewcommand{\gridpanel}[3]{%
  \begin{tikzpicture}[cam={-24}{54}{0.30}{0.30}, line cap=round, line join=round]
    \foreach \c in {0,...,15}{\foreach \r in {0,...,9}{%
      \ifcsname #1-\c-\r\endcsname \ptile{gclose}{\c}{\r}%
      \else \ptile{gfree}{\c}{\r}\fi}}
    \ptile{gstart}{1}{4}\ptile{ggoal}{14}{6}
    \draw[#3, line width=1.4pt] #2;
    \gwallsthree
  \end{tikzpicture}}
\begin{subfigure}[t]{0.32\textwidth}\centering
  \resizebox{\linewidth}{!}{\gridpanel{dj}{\dijkstraroute}{cbGreen}}
  \caption{Dijkstra's algorithm: 137 cells expanded, path cost 19.}\label{fig:grid-dijkstra}
\end{subfigure}\hfill
\begin{subfigure}[t]{0.32\textwidth}\centering
  \resizebox{\linewidth}{!}{\gridpanel{gr}{\greedyroute}{cbOrange}}
  \caption{Greedy best-first search: 48 cells expanded, path cost 35.}\label{fig:grid-greedy}
\end{subfigure}\hfill
\begin{subfigure}[t]{0.32\textwidth}\centering
  \resizebox{\linewidth}{!}{\gridpanel{as}{\astarroute}{cbGreen}}
  \caption{A* search: 81 cells expanded, path cost 19.}\label{fig:grid-astar}
\end{subfigure}

  \caption{The three algorithms on the grid. Dark tiles were expanded, and the
    line is the returned path.}\label{fig:grid-panels}
\end{figure}

Dijkstra's algorithm expands almost the whole grid. Greedy best-first search
fills the inside of the U and returns a path 16 moves longer than necessary.

A* is also drawn into the U, because the Manhattan distance ignores the wall.
Every cell inside the U is reached by moves toward the destination, so its key
equals $h(s) = 15 < 19 = C^{*}$, and Proposition~\ref{prop:expanded} forces A*
to expand it. A* then goes around the wall and returns a shortest path, after
81 expansions, all among the 137 of Dijkstra's algorithm
(Corollary~\ref{cor:subset}).

Figure~\ref{fig:timeview} raises every cell to the step at which it was
expanded: Dijkstra's algorithm spreads evenly from $s$, whereas A* first fills
the U and then climbs a narrow staircase around the wall.

\begin{figure}[htbp]
  \centering
  \begin{tabular}{@{}c@{\hspace{1em}}c@{}}
  {\footnotesize\textbf{Dijkstra}: 137 expansions} &
  {\footnotesize\textcolor{cbGreen}{\textbf{A*}}: 81 expansions, same scale}\\[2pt]
  \timeview{dj} & \timeview{as}
\end{tabular}

  \caption{Every expanded cell raised to the step at which it was expanded, on
    the same scale in both panels.}\label{fig:timeview}
\end{figure}
